\section{Source scope and verification}
% BEGIN F1.Scope
Theorem~\ref{F1.Structure} is the structural result
\cite[Theorem 1.1]{HFLQ}. Theorem~\ref{F1.Abelian} recovers its
Theorem 1.2 with the explicit marking correction. Theorem~\ref{F1.Construction}
and the exact presentation recover Theorems 1.3 and 5.1. The auxiliary
rank-three criterion, character decomposition, binomial identities,
automorphisms and maximal-class assertion are included. The original
small-group numerical census, survey discussion and general Hughes
classification are outside this principal-results contract.

The mathematical results retain their original attribution. The common
inverted-orbit and two-character arguments reorganize the source proofs;
the cyclotomic module supplies a uniform replacement for the construction
and its presentation checks. We make no novelty claim for the underlying
cyclotomic model. In the tightness converse we retain exactly
\cite[Theorem 1.1(b),(c)]{CP} in the stated odd/even range. That
primary theorem and the polyhedron correspondence were read in the
earlier accepted audit, but its full PDF is not retained locally.
There is no local byte binding or known complete-text cost for that
external proof. It is not silently replaced by finite computations.

The complete accompanying validation appendix retains three programs:
the mixed-modulus model generator, the finite marked-group checker,
and the independently authored literal-presentation checker. Together
they test seven cyclotomic and four abelian marked models and five
literal source presentations. Every tested marking is checked against
all 64 subset-intersection identities. The programs also reject the
incorrect literal marking, a corrupted cyclotomic relation and altered
source generator orders; a composite-index example shows that the
prime-index hypothesis is necessary. These are bounded checks of the
uniform proofs above, whose conclusions do not depend on a finite
enumeration or an unreported program output. The three complete
programs remain visible and chargeable as validation; both complete
historical operational drivers are inventoried separately.

The five reused Lean declarations verify the two prime-index interfaces
and three group identities. Sylow normality, the two-character theorem,
cyclotomic lattice calculations, all-parameter construction and
source-presentation isomorphism remain written proofs. The retained
tight-polyhedron classification and classical correspondence are also
outside that formal scope. This edition neither solves a new open
question nor changes the previously admitted source count. It asserts
no compression ratio; the common paper, complete application, validation
programs and external inputs have separate, explicit boundaries.
% END F1.Scope
